\documentclass[11pt,a4paper]{article}

% --- Encoding ---
\usepackage[utf8]{inputenc}
\usepackage[T1]{fontenc}
\usepackage[spanish,es-noshorthands]{babel}
\usepackage{lmodern}

% --- Layout ---
\usepackage[margin=2.5cm]{geometry}
\usepackage{parskip}
\setlength{\parindent}{0pt}
\setlength{\parskip}{0.5em plus 0.2em minus 0.1em}

% --- Tables and graphics ---
\usepackage{booktabs}
\usepackage{array}
\usepackage{enumitem}
\usepackage{xcolor}
\usepackage{graphicx}
\usepackage{tikz}
\usetikzlibrary{positioning,shapes.geometric,arrows.meta}
\usepackage{caption}

% --- Code listings ---
\usepackage{listings}
\definecolor{codegray}{rgb}{0.96,0.96,0.96}
\definecolor{codeborder}{rgb}{0.85,0.85,0.85}
\definecolor{codecomment}{rgb}{0.30,0.55,0.30}
\definecolor{codekeyword}{rgb}{0.10,0.30,0.70}
\definecolor{codestring}{rgb}{0.65,0.15,0.10}

\lstdefinestyle{cstyle}{
    language=C,
    backgroundcolor=\color{codegray},
    basicstyle=\ttfamily\footnotesize,
    breaklines=true,
    frame=single,
    rulecolor=\color{codeborder},
    showstringspaces=false,
    keywordstyle=\color{codekeyword}\bfseries,
    commentstyle=\color{codecomment}\itshape,
    stringstyle=\color{codestring},
    numbers=left,
    numberstyle=\tiny\color{gray},
    numbersep=8pt,
    aboveskip=1em,
    belowskip=1em,
}

\lstdefinestyle{shellstyle}{
    language=bash,
    backgroundcolor=\color{codegray},
    basicstyle=\ttfamily\small,
    breaklines=true,
    frame=single,
    rulecolor=\color{codeborder},
    showstringspaces=false,
    keywordstyle=\color{codekeyword}\bfseries,
    commentstyle=\color{codecomment}\itshape,
    stringstyle=\color{codestring},
    aboveskip=1em,
    belowskip=1em,
    columns=fullflexible,
}

\lstdefinestyle{asmstyle}{
    backgroundcolor=\color{codegray},
    basicstyle=\ttfamily\footnotesize,
    breaklines=true,
    frame=single,
    rulecolor=\color{codeborder},
    showstringspaces=false,
    commentstyle=\color{codecomment}\itshape,
    aboveskip=1em,
    belowskip=1em,
    columns=fullflexible,
}

\lstset{style=shellstyle}

% --- Definition box ---
\usepackage[most]{tcolorbox}
\newtcolorbox{definicion}[1][]{
    colback=blue!5!white,
    colframe=blue!50!black,
    fonttitle=\bfseries,
    title=Definición,
    sharp corners=southwest,
    rounded corners=northwest,
    arc=2pt,
    boxrule=0.6pt,
    left=8pt,right=8pt,top=4pt,bottom=4pt,
    #1
}
\newtcolorbox{idea}[1][]{
    colback=orange!5!white,
    colframe=orange!60!black,
    fonttitle=\bfseries,
    title=Idea intuitiva,
    sharp corners=southwest,
    rounded corners=northwest,
    arc=2pt,
    boxrule=0.6pt,
    left=8pt,right=8pt,top=4pt,bottom=4pt,
    #1
}

% --- Hyperlinks ---
\usepackage[colorlinks=true,
            linkcolor=blue!60!black,
            urlcolor=blue!50!black,
            citecolor=blue!60!black]{hyperref}

% --- Title ---
\title{
    \vspace{-2.5em}
    \textbf{Resumen teórico --- Sistemas Operativos} \\[0.3em]
    \large Capítulo 4 de las Notas del curso: \emph{El lenguaje C y herramientas}
}
\author{
    Tomás Rodeghiero \\
    \small Sistemas Operativos --- Universidad Nacional de Río Cuarto
}
\date{2026}

\begin{document}

\maketitle

\begin{abstract}
\noindent
Este documento es un resumen teórico del capítulo 4 de las
\emph{Notas del curso} de Sistemas Operativos (Marcelo Arroyo,
UNRC), dedicado al lenguaje C y al \emph{toolchain} de
desarrollo: pre-procesamiento, compilación, ensamblado, linking,
archivos objeto, ejecutables, símbolos, bibliotecas estáticas y
dinámicas, carga dinámica explícita y sistemas de \emph{build} como
GNU make. El objetivo es servir como material de estudio: las
definiciones formales se acompañan de explicaciones intuitivas,
ejemplos y conexiones entre conceptos, de modo que el lector pueda
preparar tanto el examen teórico como el Práctico 2 con una
comprensión sólida y articulada del proceso que transforma un
archivo \texttt{.c} en un ejecutable.
\end{abstract}

\tableofcontents
\bigskip

% =====================================================================
\section{Introducción}
% =====================================================================

Programar en C parece resumirse a un único comando ---\texttt{gcc -o
myprog main.c hello.c}---, pero detrás de él hay una secuencia de
herramientas trabajando en conjunto. Entender qué hace cada una es
clave porque el sistema operativo \emph{ejecuta} ese resultado: el
formato del binario, la disposición del código y los datos en
memoria, y la forma en que las funciones de biblioteca se enlazan en
tiempo de ejecución dependen directamente de las decisiones que se
toman en cada etapa del proceso de compilación.

\begin{definicion}
La \textbf{cadena de herramientas} (\emph{toolchain}) es el conjunto
de programas que transforman un programa fuente en un ejecutable.
Incluye, como mínimo, un \emph{pre-procesador}, un \emph{compilador},
un \emph{ensamblador} y un \emph{linker}.
\end{definicion}

\begin{idea}
Pensemos en la analogía con una receta: el código fuente es la receta
escrita en lenguaje humano, el assembly es la lista de pasos
detallados, el archivo objeto es el plato hecho pero todavía sin
montar, y el ejecutable es la comida lista para servir. El
\emph{linker} es el que arma el plato uniendo los ingredientes
(funciones) que provienen de varios módulos.
\end{idea}

% =====================================================================
\section{El lenguaje C}
% =====================================================================

C es un lenguaje de bajo nivel, diseñado por Dennis Ritchie en los
años 70, con abstracciones simples y directas al hardware. UNIX fue el
primer sistema operativo escrito casi enteramente en C, y desde
entonces el lenguaje quedó indisolublemente ligado a la programación
de sistemas: drivers, bibliotecas estándar, kernels, intérpretes y
compiladores se escriben típicamente en C.

\subsection{Tipos de datos}

\paragraph{Básicos.}
\begin{itemize}[leftmargin=*]
    \item \textbf{Enteros:} \texttt{char} (un byte), \texttt{short
        int}, \texttt{int}, \texttt{long int}. Pueden ser
        \texttt{signed} (con signo) o \texttt{unsigned} (sin signo).
        Los literales se escriben en decimal (\texttt{456}), octal
        (\texttt{0367}) o hexadecimal (\texttt{0x20AF}).
    \item \textbf{Punto flotante:} \texttt{float} (simple precisión)
        y \texttt{double} (doble precisión).
    \item \textbf{Enumeraciones:} dan nombres simbólicos a constantes
        enteras consecutivas.
\end{itemize}

\paragraph{Compuestos.}
\begin{itemize}[leftmargin=*]
    \item \textbf{Arreglos:} elementos contiguos en memoria,
        indexados desde \texttt{0} hasta \texttt{N-1}. El nombre del
        arreglo, en la mayoría de los contextos, decae a un puntero a
        su primer elemento.
    \item \textbf{Strings:} arreglos de caracteres \emph{terminados
        en} \texttt{\textbackslash 0} (\texttt{NULL}). El literal
        \texttt{"hola"} se almacena como
        \texttt{\{'h','o','l','a','\textbackslash 0'\}}.
    \item \textbf{Estructuras (\texttt{struct}):} agrupan campos
        heterogéneos bajo un mismo nombre. Cada campo se accede con
        \texttt{.}\,o, si se tiene un puntero, con \texttt{->}.
    \item \textbf{Uniones (\texttt{union}):} campos de tipos
        distintos que comparten la misma posición de memoria. Sirven
        para implementar \emph{polimorfismo} de bajo nivel
        (representaciones alternativas del mismo dato).
\end{itemize}

\paragraph{Constantes y alias.}
Los valores constantes se declaran anteponiendo la palabra clave
\texttt{const} al tipo. Los alias de tipo se introducen con
\texttt{typedef}, lo que permite escribir, por ejemplo,
\texttt{typedef unsigned char byte;} y luego usar \texttt{byte} como
sinónimo de \texttt{unsigned char}.

\subsection{Punteros}

\begin{definicion}
Un \textbf{puntero} es una variable cuyo valor es la dirección de
memoria de otra variable. Un puntero a un valor de tipo \texttt{T} se
declara como \texttt{T *p}. Tiene asociados dos operadores:
\texttt{\&v} (\emph{tomar la dirección de} \texttt{v}) y
\texttt{*p} (\emph{desreferenciar} \texttt{p}, es decir, leer el
valor apuntado).
\end{definicion}

\begin{lstlisting}[style=cstyle]
int v = 10;
int *ptr = &v;          /* ptr apunta a v */
printf("v = %d\n", *ptr);   /* desreferencia: imprime 10 */
\end{lstlisting}

En la declaración \texttt{T v[N]}, el identificador \texttt{v} es un
\emph{puntero constante} a su primer elemento; por lo tanto, \texttt{*v
== v[0]}.

\paragraph{Aritmética de punteros.}
La aritmética de punteros tiene en cuenta el tamaño del tipo apuntado.
\texttt{ptr + 1} no significa "sumar un byte" sino "avanzar al siguiente
elemento del arreglo", es decir, \texttt{ptr + sizeof(*ptr)} en bytes.

\paragraph{Punteros a función.}
C permite tomar la dirección de una función y guardarla en un puntero.
Esto habilita \emph{programación de alto orden}: pasar funciones como
parámetros, instalar callbacks, despachar dinámicamente. La sintaxis es
incómoda pero idiomática:

\begin{lstlisting}[style=cstyle]
typedef int (*fptr_t)(int, int);   /* puntero a función int x int -> int */
int inc(int v, int step) { return v + step; }
fptr_t f = inc;
printf("%d\n", f(1, 2));           /* imprime 3 */
\end{lstlisting}

\subsection{Funciones, definiciones y declaraciones}

C distingue dos conceptos relacionados pero distintos:

\begin{definicion}
La \textbf{definición} de una función incluye su \emph{cuerpo}: el
bloque de sentencias que ejecuta. Solo puede haber una definición por
función en todo el programa.
\end{definicion}

\begin{definicion}
Una \textbf{declaración} (también llamada \emph{prototipo}) especifica
únicamente el \emph{perfil} de la función ---tipo de retorno, nombre y
tipos de los parámetros--- pero no su cuerpo. Sirve para que el
compilador valide el uso cuando la definición está en otro lugar (más
abajo en el archivo o en otro módulo).
\end{definicion}

Los parámetros se pasan \emph{por valor}; el pasaje por referencia se
\emph{simula} pasando un puntero. C admite además \emph{funciones
variádicas} (con un número arbitrario de argumentos), como la propia
\texttt{printf}, cuyo perfil es
\texttt{int printf(const char *fmt, ...)}.

\subsection{Módulos: archivos \texttt{.c} y \texttt{.h}}

C no tiene una construcción sintáctica para módulos. Por convención
cada módulo se distribuye en dos archivos:

\begin{itemize}[leftmargin=*]
    \item Un archivo de implementación \texttt{nombre.c} con las
        \emph{definiciones} (cuerpo de funciones, variables globales).
    \item Un archivo de cabecera \texttt{nombre.h} con las
        \emph{declaraciones} de los símbolos exportados, pensado para
        ser incluido (\texttt{\#include "nombre.h"}) por otros
        módulos que lo usen.
\end{itemize}

Las cabeceras del sistema se incluyen con \texttt{<\,>} (por ejemplo,
\texttt{\#include <stdio.h>}); las del proyecto, con comillas
(\texttt{\#include "miarchivo.h"}). El compilador busca cada cabecera
en distintos directorios según el caso.

% =====================================================================
\section{El proceso de compilación}
% =====================================================================

\subsection{Cuatro etapas}

Compilar un programa C consiste de cuatro pasos coordinados por el
\emph{driver} \texttt{gcc} (o \texttt{clang}):

\begin{enumerate}[leftmargin=*]
    \item \textbf{Pre-procesamiento (\texttt{cpp}).} Procesa las
        directivas del programa fuente y \emph{expande} las macros y
        \texttt{\#include}. Su salida es un archivo temporal con
        extensión \texttt{.i}.
    \item \textbf{Compilación (\texttt{cc1}).} Traduce el código C ya
        expandido a \emph{assembly} de la arquitectura objetivo.
        Salida: \texttt{.s}.
    \item \textbf{Ensamblado (\texttt{as}).} Convierte el assembly en
        un \emph{archivo objeto binario} en el formato del SO
        (\texttt{ELF} en GNU/Linux, \texttt{Mach-O} en macOS,
        \texttt{PE/COFF} en Windows). Salida: \texttt{.o}.
    \item \textbf{Linking (\texttt{ld}).} Combina archivos objeto y
        bibliotecas, recalcula direcciones y resuelve referencias
        externas para producir el \emph{ejecutable} (o una nueva
        biblioteca).
\end{enumerate}

\begin{center}
\begin{tikzpicture}[
    node distance=0.55cm and 0.7cm,
    every node/.style={draw, rounded corners, font=\ttfamily\small,
                       minimum width=1.4cm, minimum height=0.65cm,
                       align=center},
    arr/.style={-{Latex[length=2mm]}, thick},
    label/.style={draw=none, font=\footnotesize\itshape}
]
\node (mainC)  {main.c};
\node (mainI)  [right=of mainC]  {main.i};
\node (mainS)  [right=of mainI]  {main.s};
\node (mainO)  [right=of mainS]  {main.o};
\node (myprog) [right=of mainO]  {myprog};

\draw[arr] (mainC) -- node[label, above]{cpp}      (mainI);
\draw[arr] (mainI) -- node[label, above]{cc1}      (mainS);
\draw[arr] (mainS) -- node[label, above]{as}       (mainO);
\draw[arr] (mainO) -- node[label, above]{ld}       (myprog);
\end{tikzpicture}
\end{center}

\subsection{Detener el proceso en una etapa}

\texttt{gcc} permite frenar la cadena de compilación en cualquier
etapa intermedia, lo que es muy útil para inspección y depuración:

\begin{table}[h]
\centering
\small
\begin{tabular}{l p{8cm} l}
\toprule
\textbf{Flag} & \textbf{Detiene tras} & \textbf{Salida} \\
\midrule
\texttt{-E} & Pre-procesamiento. & Archivo \texttt{.i} (texto). \\
\texttt{-S} & Compilación.       & Archivo \texttt{.s} (assembly). \\
\texttt{-c} & Ensamblado.        & Archivo \texttt{.o} (binario). \\
(sin flags)  & Linking.           & Ejecutable. \\
\bottomrule
\end{tabular}
\caption{Opciones de \texttt{gcc} para detener la compilación.}
\end{table}

\begin{lstlisting}
gcc -E hello.c -o hello.i     # solo preprocesa
gcc -S hello.c -o hello.s     # preprocesa + compila
gcc -c hello.c -o hello.o     # preprocesa + compila + ensambla
gcc    hello.c -o hello       # ciclo completo
\end{lstlisting}

\subsection{El rol del pre-procesador}

El pre-procesador no es un compilador: es un manipulador de texto. Sus
directivas más comunes son:

\begin{itemize}[leftmargin=*]
    \item \texttt{\#include}: copia el contenido de otro archivo.
    \item \texttt{\#define}: define una macro de sustitución textual.
    \item \texttt{\#ifdef}/\texttt{\#endif}: compilación condicional.
\end{itemize}

Por ejemplo, dado \texttt{\#define hi "Hello world"} y un cuerpo
\texttt{return hi;}, el pre-procesador entrega al compilador
\texttt{return "Hello world";}: la macro \texttt{hi} no llega siquiera
a la etapa de compilación, así que no aparece en la tabla de símbolos
del objeto.

\begin{idea}
La regla mental útil es: \emph{si está con \#, lo resuelve el
preprocesador antes de mirar el código C}. Esa es la razón por la que
las macros no respetan tipos ni alcance de bloques: son sustitución de
texto pura.
\end{idea}

% =====================================================================
\section{Archivos objeto}
% =====================================================================

\subsection{Definición y formato}

\begin{definicion}
Un \textbf{archivo objeto} (\texttt{.o}) es un binario en el formato
definido por la \emph{ABI} del sistema. Contiene el código y los
datos de un módulo en una representación que el linker puede combinar
con otros objetos para producir un ejecutable o una biblioteca.
\end{definicion}

Cada plataforma usa un formato distinto:

\begin{itemize}[leftmargin=*]
    \item \texttt{ELF} (\emph{Executable and Linkable Format}) en
        GNU/Linux y muchos UNIX modernos.
    \item \texttt{Mach-O} en macOS y otras plataformas Apple.
    \item \texttt{PE/COFF} (\emph{Portable Executable / Common Object
        File Format}) en Windows.
\end{itemize}

Los detalles cambian, pero la idea conceptual es la misma: un
\emph{header} con metadatos seguido de un conjunto de \emph{secciones}
con propósitos bien definidos.

\subsection{Secciones típicas}

\begin{itemize}[leftmargin=*]
    \item \textbf{\texttt{.text}}: instrucciones del programa (código
        ejecutable, solo lectura).
    \item \textbf{\texttt{.data}}: variables globales y estáticas con
        valores iniciales no nulos. Se subdivide en datos
        \emph{inicializados} y \emph{no inicializados} (esto último a
        veces aparece como \texttt{.bss}).
    \item \textbf{\texttt{.rodata}} (en macOS:
        \texttt{\_\_TEXT,\_\_cstring}): constantes y literales de
        cadena, de solo lectura.
    \item \textbf{Symbol table}: mapping
        \emph{símbolo $\to$ dirección} que el linker usa para resolver
        referencias.
    \item \textbf{Reallocation entries}: lista de instrucciones cuyos
        operandos hacen referencia a símbolos externos y necesitan ser
        \emph{parchadas} por el linker.
\end{itemize}

\subsection{Símbolos: definidos y no definidos}

\begin{definicion}
Un \textbf{símbolo} es un identificador (función, variable global,
constante) con una dirección dentro del archivo objeto. Si la
dirección está asignada, el símbolo está \emph{definido} en ese
módulo; si solo aparece como referencia sin definición, está
\emph{indefinido} (\texttt{*UND*}) y el linker deberá resolverlo
buscando el módulo o la biblioteca que lo defina.
\end{definicion}

La tabla de símbolos es un \emph{mapping} de la forma
$\textit{symtbl}(\textit{symbol}) \to \textit{address}$. Por ejemplo,
en un archivo \texttt{main.o} que llama a \texttt{printf} y a
\texttt{hello}, ambos símbolos aparecen como \texttt{*UND*} porque
sus definiciones viven en otros módulos: \texttt{libc} y
\texttt{hello.o}, respectivamente.

\subsection{Tabla de reubicación}

\begin{definicion}
La \textbf{tabla de reubicación} (\emph{relocation table}) es un
mapping de la forma
$\textit{relloc}(\textit{instr\_address}, \textit{instr\_type}) \to
\textit{external\_symbol}$ que indica al linker qué operandos de qué
instrucciones debe completar al combinar los objetos.
\end{definicion}

Las entradas de la tabla de reubicación corresponden, típicamente, a
instrucciones de la forma \texttt{call <addr>}, \texttt{load/store
<addr>}, etc., donde \texttt{<addr>} todavía no se conoce porque el
símbolo es externo. El compilador deja \texttt{0} (o un valor de
relleno) en el operando y agrega la entrada correspondiente.

\subsection{Direcciones en un objeto relocatable}

En un archivo objeto cada función empieza, típicamente, en el
\emph{offset} \texttt{0x0} de su sección \texttt{.text}: cada módulo
"se siente solo en el mundo". Las direcciones definitivas las asigna
el linker cuando combina varios objetos en un mismo binario.

% =====================================================================
\section{Enlazado (linking)}
% =====================================================================

\subsection{Función del linker}

\begin{definicion}
El \textbf{linker} (también \emph{link-editor}) es la herramienta
encargada de producir binarios (ejecutables o bibliotecas) a partir de
la \emph{composición} de varios archivos objeto y bibliotecas.
\end{definicion}

Su trabajo se descompone, conceptualmente, en dos pasos:

\begin{enumerate}[leftmargin=*]
    \item \textbf{Concatenar} las secciones de código y datos de cada
        archivo objeto y \emph{recalcular} las direcciones de
        variables, constantes y funciones. Cada \texttt{.text} que
        empezaba en offset \texttt{0} pasa a ocupar una dirección
        absoluta en el ejecutable.
    \item \textbf{Resolver las referencias externas}: completar los
        operandos de las instrucciones de llamada y de los accesos a
        variables globales externas, usando para ello las tablas de
        símbolos y de reubicación de cada objeto.
\end{enumerate}

\subsection{Resolución estática}

\begin{definicion}
El \textbf{linking estático} es la resolución de referencias externas
en \emph{tiempo de compilación}: el linker copia el código necesario
dentro del ejecutable final y completa todas las direcciones.
\end{definicion}

Si el módulo \texttt{main.o} llama a \texttt{hello} (definido en
\texttt{hello.o}), el linker, al combinar ambos, conoce la dirección
final de \texttt{\_hello} y completa el operando del \texttt{call}
con esa dirección. La invocación queda \emph{resuelta} y el
ejecutable es \emph{auto-contenido} respecto de \texttt{hello}.

\subsection{Resolución dinámica (introducción)}

\begin{definicion}
El \textbf{linking dinámico} es la resolución de referencias externas
en \emph{tiempo de carga} o de \emph{primera invocación}: el código
de la biblioteca no se copia al ejecutable, sino que se mantiene en
una biblioteca compartida que el SO carga aparte.
\end{definicion}

Las funciones de la libc, como \texttt{printf}, viven en una
biblioteca compartida (\texttt{libc.so}, \texttt{libSystem.B.dylib}).
El linker estático no copia el código de \texttt{printf} al
ejecutable: solo deja una referencia que será resuelta por el
\emph{dynamic linker} cuando el programa se cargue. Más adelante se
describe el mecanismo en detalle.

\subsection{Ejemplo completo: tres llamadas}

Ejemplificando con los programas \texttt{main.c} y \texttt{hello.c} de
las notas: \texttt{main.o} contiene dos referencias externas
(\texttt{\_hello} y \texttt{\_printf}). Antes del linking, el
disassembly muestra ambos \texttt{call} apuntando a \texttt{0} (sin
resolver):

\begin{lstlisting}[style=asmstyle]
.text section (main.o)
_main:
0x00  ...
0x18  call 0x18         ; <-- _hello, sin resolver
0x1c  call 0x1c         ; <-- _printf, sin resolver
\end{lstlisting}

Luego del linking, en el ejecutable \texttt{myprog}:

\begin{lstlisting}[style=asmstyle]
.text section (myprog)
main:
0x...   call hello              ; resuelto estaticamente
0x...   call printf_stub        ; resolverá dyld
hello:
0x...   ...
        ret
\end{lstlisting}

El \texttt{call} a \texttt{hello} apunta directamente al inicio de la
función \texttt{\_hello} dentro del ejecutable. El \texttt{call} a
\texttt{printf} apunta a un \emph{stub} (más abajo).

% =====================================================================
\section{Bibliotecas}
% =====================================================================

\begin{definicion}
Una \textbf{biblioteca} es una colección de módulos (código y datos)
\emph{reusables}. Generalmente contiene un conjunto de archivos
objeto. Pueden ser de dos tipos: para \emph{enlazado estático} o de
\emph{enlazado dinámico (shared)}.
\end{definicion}

\subsection{Bibliotecas estáticas}

\begin{definicion}
Una \textbf{biblioteca estática} es un \emph{contenedor} (formato
\emph{archiver} en sistemas tipo UNIX) de archivos objeto. Incluye un
\emph{índice} con los nombres de los símbolos exportados por cada
miembro.
\end{definicion}

En sistemas tipo UNIX se generan con la herramienta \texttt{ar}:

\begin{lstlisting}
ar rcs libmilib.a file1.o file2.o ...
\end{lstlisting}

Las opciones significan: \texttt{r} reemplaza miembros existentes,
\texttt{c} crea el archivo si no existe y \texttt{s} construye o
actualiza el índice de símbolos. Por convención, una biblioteca
\texttt{B} se nombra \texttt{libB.a}.

Cuando el linker busca un símbolo no resuelto, consulta el índice de
las bibliotecas estáticas pasadas en la línea de comandos. Si lo
encuentra, copia \emph{solo el archivo objeto necesario} dentro del
ejecutable.

\begin{lstlisting}
gcc -o myprog -L. -lhello main.o
\end{lstlisting}

El flag \texttt{-L\$DIR} agrega \texttt{\$DIR} al \emph{path} de
búsqueda; \texttt{-lhello} le dice al linker que busque
\texttt{libhello.a} (o, si está disponible, su versión dinámica).

\subsection{Bibliotecas dinámicas}

\begin{definicion}
Una \textbf{biblioteca dinámica} (también \emph{shared library}) es
un objeto compartible que el SO carga y enlaza bajo demanda en
\emph{tiempo de ejecución}. Se nombra \texttt{libfoo.so} en GNU/Linux,
\texttt{libfoo.dylib} en macOS y \texttt{foo.dll} en Windows.
\end{definicion}

\paragraph{Generación.}
Con \texttt{gcc} se generan en dos pasos: compilar a objeto con
código \emph{independiente de posición} (\texttt{-fPIC}) y luego
linkear con la opción \texttt{-shared}.

\begin{lstlisting}
gcc -c -fPIC hello.c             # genera hello.o relocatable
gcc -shared -o libhello.so hello.o
\end{lstlisting}

\paragraph{Ventajas y desventajas.}
\begin{itemize}[leftmargin=*]
    \item \textbf{Ventaja:} el ejecutable no incluye el código de la
        biblioteca; varios procesos pueden compartir una sola copia
        cargada en memoria. Las actualizaciones de la biblioteca
        benefician a todos los programas sin recompilar.
    \item \textbf{Desventaja:} sobrecarga al inicio del proceso
        (resolución de símbolos) y dependencia explícita: si la
        biblioteca no está, el programa no arranca. Se mitiga con
        \emph{lazy linking}.
\end{itemize}

\paragraph{Preferencia del linker.}
Si conviven una versión estática y una dinámica de la misma
biblioteca, el linker \emph{prefiere la dinámica} salvo que se le
indique explícitamente lo contrario (por ejemplo, pasando
\texttt{libhello.a} directamente).

\subsection{Lazy linking}

\begin{definicion}
\textbf{Lazy linking} es la técnica por la cual las funciones de una
biblioteca dinámica no se resuelven todas al cargar el programa,
sino \emph{cada una en su primera invocación}. Reduce la latencia
inicial del proceso a costo de un pequeño salto extra en cada
llamada.
\end{definicion}

\paragraph{Mecanismo: stubs y GOT.}
El linker estático genera, para cada función externa, un pequeño
fragmento de código llamado \emph{stub} (ej.\ \texttt{printf\_stub}) y
una entrada en una tabla indirecta llamada \textbf{Global Offset
Table (GOT)}. Esquemáticamente:

\begin{enumerate}[leftmargin=*]
    \item En el ejecutable, \texttt{call printf} apunta al stub.
    \item El stub realiza un \texttt{jmp} indirecto leyendo la
        dirección de \texttt{GOT[printf]}.
    \item Inicialmente, \texttt{GOT[printf]} apunta a una rutina de
        resolución del \emph{dynamic linker}.
    \item En la primera invocación, el dynamic linker localiza
        \texttt{printf} en \texttt{libc.so}, escribe su dirección
        verdadera en \texttt{GOT[printf]} y salta a ella.
    \item Las invocaciones siguientes ya encuentran la dirección
        correcta en la GOT y saltan directamente, sin volver a
        invocar al dynamic linker.
\end{enumerate}

\begin{idea}
Los \emph{stubs} usan \texttt{jmp} y no \texttt{call} porque deben
retornar al programa que invocó originalmente, no al stub. Es la
razón por la que la "indirección" no agrega un nivel innecesario en
la pila.
\end{idea}

\subsection{Carga dinámica explícita: \texttt{dlopen}/\texttt{dlsym}}

Hasta ahora se asume que las dependencias dinámicas se declaran en
tiempo de linking (\texttt{-lhello}). Existe otra opción: la
\textbf{carga dinámica explícita} mediante la API POSIX
\texttt{<dlfcn.h>}, que permite a un programa cargar una biblioteca
\emph{en tiempo de ejecución} sin haberla declarado al compilar.

\begin{itemize}[leftmargin=*]
    \item \texttt{dlopen("ruta", flags)}: carga la biblioteca y
        devuelve un \emph{handle} opaco. Los flags más usados son
        \texttt{RTLD\_NOW} (resolver todos los símbolos al cargar) y
        \texttt{RTLD\_LAZY} (resolverlos bajo demanda).
    \item \texttt{dlsym(handle, "nombre")}: resuelve el símbolo
        dentro de la biblioteca cargada y devuelve su dirección.
    \item \texttt{dlclose(handle)}: descarga la biblioteca.
    \item \texttt{dlerror()}: devuelve el último mensaje de error y
        limpia el estado interno.
\end{itemize}

\begin{lstlisting}[style=cstyle]
#include <dlfcn.h>

void *handle = dlopen("./libhello.so", RTLD_NOW);
char *(*fn)(void) = dlsym(handle, "hello");
printf("%s\n", fn());
dlclose(handle);
\end{lstlisting}

Esta técnica es la base de los sistemas de \emph{plugins} y de
cualquier mecanismo de extensibilidad en tiempo de ejecución
(módulos cargables, backends seleccionables por configuración, etc.).

% =====================================================================
\section{Espacio de direcciones del ejecutable}
% =====================================================================

Cuando el SO carga un ejecutable, el espacio de direcciones del
proceso queda dispuesto, conceptualmente, así (de direcciones bajas
a altas):

\begin{center}
\begin{tikzpicture}[
    every node/.style={draw, font=\ttfamily\small,
                       minimum width=4cm, minimum height=0.65cm},
    label/.style={draw=none, font=\footnotesize\itshape}
]
\node (text) at (0,0)        {.text  (codigo)};
\node (rodata) [above=0pt of text]  {.rodata (constantes)};
\node (data)   [above=0pt of rodata]{.data  (globales inicializadas)};
\node (bss)    [above=0pt of data]  {.bss   (globales no inicializadas)};
\node (heap)   [above=0pt of bss]   {heap (crece $\uparrow$)};
\node (gap)    [above=0pt of heap, draw=none] {\vdots};
\node (stack)  [above=0pt of gap]   {stack (crece $\downarrow$)};
\node[draw=none, right=0.2cm of stack, label]   {direcciones altas};
\node[draw=none, right=0.2cm of text, label]    {direcciones bajas};
\end{tikzpicture}
\end{center}

\begin{itemize}[leftmargin=*]
    \item \textbf{Texto (\texttt{.text})}: instrucciones del
        programa. De solo lectura, compartible entre procesos.
    \item \textbf{Datos inicializados (\texttt{.data})}: variables
        globales y estáticas con valor inicial no nulo.
    \item \textbf{Datos no inicializados (\texttt{.bss})}: variables
        globales y estáticas sin valor inicial; el kernel las inicia
        a cero.
    \item \textbf{Heap}: memoria dinámica obtenida con
        \texttt{malloc}/\texttt{brk}; crece hacia direcciones altas.
    \item \textbf{Stack}: pila de llamadas; crece hacia direcciones
        bajas. Cada invocación a una función crea un \emph{stack
        frame} con sus parámetros, variables locales, dirección de
        retorno y registros guardados.
\end{itemize}

\begin{idea}
Una variable local \texttt{int x;} dentro de una función vive en el
\emph{stack frame} de esa función y desaparece cuando la función
retorna. Una variable global \texttt{int g = 42;} vive en
\texttt{.data} durante toda la ejecución del programa. Una constante
literal como \texttt{"hola"} vive en \texttt{.rodata} y es de solo
lectura: intentar modificarla suele provocar \texttt{SIGSEGV}.
\end{idea}

% =====================================================================
\section{Build systems: \texttt{make} y los Makefiles}
% =====================================================================

\subsection{Motivación}

En un proyecto real, las cuatro etapas del toolchain se ejecutan
sobre decenas o cientos de archivos. Recompilar todo desde cero ante
cada cambio es inviable. Un \emph{build system} automatiza el proceso
y, sobre todo, evita el trabajo redundante: solo reconstruye lo que
realmente cambió.

\subsection{GNU make}

\begin{definicion}
\texttt{make} es una herramienta que toma una especificación de
\emph{targets} y \emph{dependencias} (escrita en un archivo
\texttt{Makefile}) y construye automáticamente el grafo de
dependencias. Comparando los \emph{timestamps} de los archivos,
decide qué reglas disparar y en qué orden.
\end{definicion}

La forma básica de una regla es:

\begin{lstlisting}[style=asmstyle]
target: dependencias
<TAB>comando
\end{lstlisting}

\textbf{Importante:} la indentación debe ser exactamente un
\texttt{TAB}, no espacios.

Ejemplo idiomático para construir \texttt{myprog} a partir de
\texttt{main.c} y \texttt{hello.c}:

\begin{lstlisting}[style=asmstyle]
myprog: main.o hello.o
	gcc -o myprog main.o hello.o

main.o: main.c
	gcc -c main.c

hello.o: hello.c
	gcc -c hello.c
\end{lstlisting}

\subsection{Variables automáticas y reglas de patrón}

Para no repetir información, \texttt{make} ofrece \emph{variables
automáticas}:

\begin{itemize}[leftmargin=*]
    \item \texttt{\$@}: nombre del target.
    \item \texttt{\$\^{}}: lista completa de dependencias.
    \item \texttt{\$<}: primera dependencia (la "fuente").
\end{itemize}

Y \emph{reglas de patrón} que generalizan la construcción:

\begin{lstlisting}[style=asmstyle]
myprog: main.o hello.o
	gcc -o $@ $^

%.o: %.c
	gcc -c $<
\end{lstlisting}

\subsection{Compilación incremental}

\texttt{make} construye un grafo de dependencias y, en cada
invocación, decide:

\begin{itemize}[leftmargin=*]
    \item Si el target no existe o es más viejo que alguna
        dependencia, hay que reconstruirlo.
    \item Antes de hacerlo, hay que asegurarse de que las
        dependencias estén actualizadas (recursión).
\end{itemize}

\begin{idea}
Si \texttt{main.c} no cambió, \texttt{make} no recompila
\texttt{main.o}, aunque sí pueda regenerar \texttt{hello.o} y
relinkear el ejecutable. Es la diferencia entre un proyecto
manejable y uno inviable.
\end{idea}

\subsection{Autotools y otros sistemas de build}

Sobre \texttt{make} se construye el ecosistema de las
\textbf{autotools}:

\begin{itemize}[leftmargin=*]
    \item \textbf{Autoconf:} genera el script \texttt{configure} a
        partir de un archivo \texttt{configure.ac}; el script detecta
        el entorno (compilador disponible, opciones) y produce un
        \texttt{Makefile} adaptado.
    \item \textbf{Automake:} a partir de \texttt{Makefile.am} y
        \texttt{configure.ac} genera el \texttt{Makefile.in} portable
        que usará \texttt{configure}.
\end{itemize}

El flujo típico para un usuario que recibe el \emph{tarball} de un
proyecto es:

\begin{lstlisting}
tar xzvf my-package.tgz
cd my-package
./configure && make && sudo make install
\end{lstlisting}

Otros sistemas modernos (\texttt{CMake}, \texttt{Meson}, \texttt{Ninja},
\texttt{Bazel}) parten de las mismas ideas y agregan portabilidad,
paralelismo o reproducibilidad.

% =====================================================================
\section{Resumen integrador}
% =====================================================================

Conviene cerrar visualizando todo el flujo en un único ejemplo. Sea el
comando

\begin{lstlisting}
gcc -o myprog main.c hello.c
\end{lstlisting}

Lo que ocurre detrás es lo siguiente:

\begin{enumerate}[leftmargin=*]
    \item \texttt{gcc} invoca al \emph{pre-procesador} sobre cada
        \texttt{.c}: expande \texttt{\#include}, resuelve
        \texttt{\#define} y produce un \texttt{.i} temporal.
    \item Pasa el \texttt{.i} al \emph{compilador} (\texttt{cc1}),
        que lo traduce a \emph{assembly} específico de la
        arquitectura (\texttt{.s}).
    \item Invoca al \emph{ensamblador} (\texttt{as}) sobre cada
        \texttt{.s} y obtiene los archivos objeto \texttt{main.o} y
        \texttt{hello.o}, en formato ELF/Mach-O. En cada uno, las
        funciones empiezan en offset \texttt{0x0} y las llamadas a
        símbolos externos quedan sin resolver, registradas en la
        tabla de reubicación.
    \item Llama al \emph{linker} (\texttt{ld}), que combina los
        objetos: concatena las secciones, recalcula direcciones,
        resuelve estáticamente \texttt{\_hello} (al estar definido
        en \texttt{hello.o}) y deja \texttt{\_printf} pendiente, con
        un \emph{stub} y una entrada en la GOT, para que lo resuelva
        el \emph{dynamic linker} en ejecución.
    \item El binario resultante (\texttt{myprog}) tiene una sección
        \texttt{.interp} que indica el dynamic linker que debe
        cargarlo (\texttt{/lib/ld-linux.so} en Linux,
        \texttt{/usr/lib/dyld} en macOS) y una sección
        \texttt{.dynamic} con la lista de bibliotecas dinámicas que
        necesita (la libc, principalmente).
    \item Al ejecutar \texttt{./myprog}, el SO carga el ejecutable,
        invoca al dynamic linker, este carga \texttt{libc} (si no
        está ya en memoria), inicializa la GOT y, finalmente, el
        kernel salta a \texttt{\_start} (que termina invocando a
        \texttt{main}).
    \item Cuando \texttt{main} llama por primera vez a
        \texttt{printf}, el stub salta vía la GOT al dynamic linker,
        que resuelve la dirección, la escribe en la GOT y salta a la
        verdadera \texttt{printf} de la libc. La siguiente invocación
        ya va directa.
\end{enumerate}

Detrás de un comando aparentemente simple operan, en conjunto: el
\emph{pre-procesador}, el \emph{compilador}, el \emph{ensamblador},
el \emph{linker estático}, el \emph{loader} del SO, el
\emph{dynamic linker}, las \emph{bibliotecas compartidas} y la
\emph{tabla de reubicación} de cada objeto. El estudio de la
compilación y el linking consiste, precisamente, en entender cómo
todas esas piezas funcionan juntas para transformar texto en código
en ejecución.

% =====================================================================
\section*{Referencias}
\addcontentsline{toc}{section}{Referencias}
% =====================================================================

\begin{enumerate}[leftmargin=*]
    \item Marcelo Arroyo. \emph{Notas del curso de Sistemas
        Operativos}, capítulo \emph{El lenguaje C y herramientas}.
        UNRC, 2026.\\
        \url{https://marceloarroyo.gitlab.io/cursos/SO/index.html\#/toolchain}
    \item Brian Kernighan, Dennis Ritchie. \emph{The C programming
        language}, 2.\textsuperscript{a} edición. Prentice Hall, 1988.
    \item John R.\ Levine. \emph{Linkers and Loaders}. Morgan
        Kaufmann, 2000.
    \item TIS Committee. \emph{Tool Interface Standard (TIS)
        Executable and Linking Format (ELF) Specification}, v1.2.
    \item POSIX.1-2017 / Single UNIX Specification. The Open Group,
        2018.
    \item Páginas de manual: \texttt{man 1 gcc}, \texttt{man 1 ld},
        \texttt{man 1 ar}, \texttt{man 1 objdump}, \texttt{man 1 nm},
        \texttt{man 3 dlopen}, \texttt{man 1 make}.
\end{enumerate}

\end{document}
